\documentclass{article}
\usepackage{graphicx} % Required for inserting images

\title{isolation forest}
\author{Shuvasish Saha}
\date{April 2026}

\begin{document}

\maketitle
\subsubsection{The Anomaly Score Formula}
The anomaly score $s(x, n)$ for an instance $x$ is defined as:
\begin{equation}
    s(x, n) = 2^{-\frac{E(h(x))}{c(n)}}
\end{equation}

Where $c(n)$ is the normalization factor calculated as:
\begin{equation}
    c(n) = 2 \ln(n - 1) + 0.5772156649 - \frac{2(n - 1)}{n}
\end{equation}

\textbf{Variable Breakdown:}
\begin{itemize}
    \item \textbf{$x$}: The specific smart contract instance being analyzed.
    \item \textbf{$n$}: The number of instances used to build the isolation trees.
    \item \textbf{$h(x)$}: The path length (number of edges) from the root node to the external leaf node where $x$ terminates.
    \item \textbf{$E(h(x))$}: The Expected (Average) Path Length of $x$ across a collection of isolation trees (iTrees).
\end{itemize}
\subsubsection{Autoencoder-Based Detection}
The Autoencoder identifies anomalies by calculating the reconstruction error for each instance. The anomaly score is defined using the Mean Squared Error (MSE):

\begin{equation}
    L(x, \hat{x}) = \frac{1}{m} \sum_{j=1}^{m} (x_j - \hat{x}_j)^2
\end{equation}

\textbf{Variable Breakdown:}
\begin{itemize}
    \item \textbf{$x$}: The original input feature vector representing the smart contract.
    \item \textbf{$\hat{x}$}: The reconstructed feature vector produced by the decoder.
    \item \textbf{$m$}: The total number of input features in the dataset.
    \item \textbf{$L(x, \hat{x})$}: The reconstruction loss, which serves as the indicator for potential fraudulent activity.
\end{itemize}
\subsubsection{XGBoost Objective Function and Regularization}
The final classification layer employs XGBoost, which minimizes a regularized objective function at iteration $t$:

\begin{equation}
    \text{Obj}^{(t)} = \sum_{i=1}^{n} L(y_i, \hat{y}_i^{(t)}) + \sum_{k=1}^{t} \Omega(f_k)
\end{equation}

To mitigate overfitting, the complexity of each tree $f$ is constrained by the regularization term $\Omega(f)$, incorporating $L1$ ($\alpha$) and $L2$ ($\lambda$) penalties:

\begin{equation}
    \Omega(f) = \gamma T + \frac{1}{2} \lambda \|w\|^2 + \alpha \|w\|
\end{equation}

\textbf{Variable Breakdown:}
\begin{itemize}
    \item \textbf{$L$}: Differentiable loss function measuring prediction error.
    \item \textbf{$T$}: Total number of leaf nodes in the tree.
    \item \textbf{$w$}: The leaf weight vector.
    \item \textbf{$\gamma$}: The complexity cost for adding an additional leaf.
    \item \textbf{$\lambda, \alpha$}: Regularization coefficients for $L2$ and $L1$ norms, respectively.
\end{itemize}

\section{Introduction}

\end{document}
